\subsection{Study 1 Results}
We first summarize the analysis sample. After dropping test IDs, 337 participant CSV files remained. Among them, 335 participants had complete treatment data across all three interaction-availability groups (Static, Interactive, and Buggy), yielding 1{,}005 rows for each hypothesis analysis. After normalization, all trust and interaction questionnaire items fell within the expected 1--7 range.

\noindent\textbf{Results: Confidence Shift from Baseline by Visualization Family.}
We first examined whether the uncertainty visualizations changed confidence relative to each participant's baseline forecast judgment. A linear model with participant fixed effects showed a significant effect of visualization family across Baseline, PI, Ensemble, and PI+Ensemble, $F(3, 779)=15.90$, $p<.001$. Estimated marginal means were Baseline ($M=4.40$, $SE=0.05$), PI ($M=4.19$, $SE=0.06$), Ensemble ($M=4.11$, $SE=0.06$), and PI+Ensemble ($M=3.96$, $SE=0.05$). Relative to baseline, all three uncertainty visualizations reduced confidence on average (PI: $M_{\Delta}=-0.21$, $SE=0.07$, $n=226$; Ensemble: $M_{\Delta}=-0.31$, $SE=0.09$, $n=221$; PI+Ensemble: $M_{\Delta}=-0.44$, $SE=0.07$, $n=335$).

\noindent\textbf{Takeaway.}
Introducing uncertainty information lowered confidence relative to baseline judgments, and the combined PI+Ensemble design produced the largest reduction.

\noindent\textbf{Results: Trust Composite Across Interaction Availability and Visualization Family.}
We next summarized trust ratings on the 1--7 composite scale. Across interaction availability, estimated trust was highest for Interactive ($M=4.73$, $SE=0.05$), followed by Buggy ($M=4.42$, $SE=0.05$), then Static ($M=4.34$, $SE=0.05$). Across visualization families, PI had the highest trust ($M=4.82$, $SE=0.05$), followed by Ensemble ($M=4.37$, $SE=0.06$) and PI+Ensemble ($M=4.29$, $SE=0.05$). At the condition level, Interactive+PI was highest ($M=4.97$, $SE=0.09$), while Static+Ensemble ($M=4.05$, $SE=0.10$) and Static+PI+Ensemble ($M=4.05$, $SE=0.10$) were lowest.

\noindent\textbf{Takeaway.}
Working interaction was associated with higher trust than both static and buggy versions, and PI was the most trusted visualization family overall.

\noindent\textbf{Results: Interaction Experience Across Interaction Availability.}
We then examined the interaction experience composite. Interactive again showed the highest estimated mean ($M=3.86$, $SE=0.03$), followed by Buggy ($M=3.73$, $SE=0.03$) and Static ($M=3.44$, $SE=0.03$). Across visualization families, PI was highest ($M=3.77$, $SE=0.04$), followed by Ensemble ($M=3.68$, $SE=0.04$) and PI+Ensemble ($M=3.58$, $SE=0.03$). Condition-level estimates ranged from Interactive+Ensemble ($M=3.95$, $SE=0.07$) to Static+PI+Ensemble ($M=3.31$, $SE=0.07$).

\noindent\textbf{Takeaway.}
Direct interaction improved the subjective interaction experience, while buggy interaction reduced that benefit but still performed better than static displays on average.

\noindent\textbf{Results: Visualization Family Within Interactive Availability Only.}
Because trust differences were most pronounced under Interactive availability, we ran a follow-up one-way ANOVA within the Interactive rows only. Visualization family significantly affected trust, $F(2, 332)=5.11$, $p=.006$. Means were highest for PI ($n=117$, $M=4.99$, $SD=1.09$, $SE=0.10$), followed by Ensemble ($n=109$, $M=4.73$, $SD=1.11$, $SE=0.11$), and PI+Ensemble ($n=109$, $M=4.48$, $SD=1.33$, $SE=0.13$).

\noindent\textbf{Takeaway.}
Among the working interactive conditions, PI elicited the highest trust.
